\clearpage
\VerdierPhysicalPageAnchor{0120}%
\begingroup
\small
\setlength{\abovedisplayskip}{4pt plus 2pt minus 2pt}%
\setlength{\belowdisplayskip}{4pt plus 2pt minus 2pt}%

\VerdierSemanticAnchor{ch2-1-2-9-opposite-category}
$\mathrm{ii)}\Longrightarrow\mathrm{iii)}$ pour la catégorie opposée à
$\mathcal D$. L'implication $\mathrm{i)}\Longrightarrow\mathrm{iii)}$ pour
toute catégorie triangulée est donc équivalente à l'implication
$\mathrm{ii)}\Longrightarrow\mathrm{iii)}$ pour toute catégorie triangulée.
Pour démontrer les implications :

\VerdierSemanticAnchor{ch2-1-2-9-reduction-implication}
\[
  \mathrm{i)}\Longleftrightarrow\mathrm{ii)}
  \Longleftrightarrow\mathrm{iii)}\Longrightarrow\mathrm{iv)}\,,
\]
il suffit donc de démontrer l'implication
$\mathrm{i)}\Longrightarrow\mathrm{iii)}$.

\VerdierSemanticAnchor{ch2-1-2-10}
\medskip
\noindent\textbf{1.2.10. Démonstration de la proposition 1.2.9. (suite).}
Démontrons d'abord qu'un monomorphisme est décomposable. Soient
$f\colon X\longrightarrow Y$ un monomorphisme et
$X\xrightarrow{f}Y\xrightarrow{v}Z\xrightarrow{w}X[1]$ un triangle
distingué contenant $f$ (TRI).

\VerdierSemanticAnchor{ch2-1-2-10-monomorphism}
Le morphisme $f$ étant un monomorphisme, il en est de même de $f[1]$. Il
résulte alors de TRII et de (1.2.2) que $w=0$. Il résulte alors de (1.2.6)
que le morphisme $f$ est décomposable.

\VerdierSemanticAnchor{ch2-1-2-10-kernel-decomposition}
Démontrons maintenant que
$\mathrm{i)}\Longrightarrow\mathrm{iii)}$ dans le cas général. Supposons
que le morphisme $f\colon X\longrightarrow Y$ admette un noyau
$\operatorname{Ker}(f)\longrightarrow X$. Le morphisme
$\operatorname{Ker}(f)\longrightarrow X$ est un monomorphisme, donc il est
décomposable. Le morphisme $f$ est donc isomorphe à un morphisme du type :

\VerdierSemanticAnchor{ch2-1-2-10-kernel-normal-form}
\[
  \operatorname{Ker}(f)\oplus X'
  \xrightarrow{(0,g)}Y\,,
\]
dont le noyau est

\VerdierSemanticAnchor{ch2-1-2-10-kernel-inclusion}
\[
  \operatorname{Ker}(f)
  \xrightarrow{\left(\begin{smallmatrix}\operatorname{id}\\0\end{smallmatrix}\right)}
  \operatorname{Ker}(f)\oplus X'.
\]
On en déduit immédiatement que le morphisme $g$ est un monomorphisme; par
suite, il est décomposable. Le morphisme $f\colon X\longrightarrow Y$ est
donc isomorphe à un morphisme du type :

\VerdierSemanticAnchor{ch2-1-2-10-final-matrix}
\[
  \operatorname{Ker}(f)\oplus X'
  \xrightarrow{\left(\begin{smallmatrix}0&\operatorname{id}\\0&0\end{smallmatrix}\right)}
  X'\oplus Y'.
\]
Il est donc décomposable.

\VerdierSemanticAnchor{ch2-1-2-11}
\medskip
\noindent\textbf{1.2.11. Démonstration de la proposition 1.2.9.
(fin)\footnote{Nous utilisons ici un argument dû à Freyd.}}
Il reste à démontrer l'implication
$\mathrm{iv)}\Longrightarrow\mathrm{iii)}$ sous les hypothèses
supplémentaires.

\VerdierSemanticAnchor{ch2-1-2-11-duality-reduction}
Remarquons tout d'abord que les sommes directes dénombrables sont
représentables dans $\mathcal D$ si et seulement si les produits directs
dénombrables sont représentables dans $\mathcal D^{\circ}$. De plus,
l'implication $\mathrm{iv)}\Longrightarrow\mathrm{iii)}$ pour la catégorie
$\mathcal D$ est équivalente à l'implication
$\mathrm{iv)}\Longrightarrow\mathrm{iii)}$ pour la catégorie opposée à
$\mathcal D$.

\VerdierSemanticAnchor{ch2-1-2-11-product-case-boundary}
Il suffit donc de démontrer l'implication
$\mathrm{iv)}\Longrightarrow\mathrm{iii)}$ lorsque les produits
dénombrables sont représentables dans $\mathcal D$, \textit{i.e.} compte
tenu de ce qui précède, de
\endgroup
